% Créditos de las fotografías y de las ilustraciones generadas por IA del
% libro 4 (Química universitaria, tercer año), edición española. Los
% registros de licencia legibles por máquina están en
% images/book4/CREDITS.md; mantener ambos archivos sincronizados con la
% versión inglesa (frontmatter/image-credits-book4.tex).
\chapter*{Créditos de las imágenes}

Los dibujos de este libro son originales. Las fotografías, cuando se
usan, se emplean bajo sus respectivas licencias libres, con nuestro
agradecimiento a sus autores; los enlaces completos a las fuentes y las
URL de las licencias de cada fotografía figuran en
\texttt{images/book4/CREDITS.md}, en el repositorio del libro.
Los pictogramas de peligro son los del Sistema Globalmente Armonizado de
clasificación y etiquetado de productos químicos, tal como los dibuja el
paquete \texttt{ghsystem} de \TeX\ Live.

\bigskip
Junto a las fotografías y a los dibujos originales, algunas figuras son
ilustraciones generadas por IA, producidas con la generación de imágenes
de OpenAI a partir de instrucciones escritas por los editores del libro,
y revisadas en cuanto a su exactitud química antes de incluirlas. La
instrucción completa de cada imagen generada está registrada en
\texttt{images/book4/ai/PROMPTS.md}, en el repositorio.
\bigskip
\noindent Fotografías, por capítulo:
\begin{itemize}\small
  \item Cap.~1: Erwin Schr\"odinger, 1933, Fundación Nobel, dominio público
        (Wikimedia Commons).
  \item Cap.~4: cristales de nieve, Wilson A. Bentley (hacia 1902), dominio público
        (Wikimedia Commons).
  \item Cap.~6: radioantenas en una meseta de gran altitud, ESO/B.~Tafreshi
        (twanight.org), CC BY 4.0 (Wikimedia Commons).
  \item Cap.~7: fluorita a la luz del día y bajo luz ultravioleta, de Masha
        Milshina, CC BY 4.0 (Wikimedia Commons).
  \item Cap.~8: un laboratorio de RMN, de Shandchem, CC BY 2.0 (Wikimedia
        Commons).
  \item Cap.~9: William Lawrence Bragg, Fundación Nobel, dominio público
        (Wikimedia Commons).
  \item Cap.~10: la tumba de Ludwig Boltzmann, de Feldkurat Katz,
        CC BY-SA 4.0 (Wikimedia Commons).
  \item Cap.~13: espirales de Belousov--Zhabotinsky, de Stephen Morris, CC BY 2.0
        (Wikimedia Commons).
  \item Cap.~14: el agujero de ozono antártico, NASA, dominio público (Wikimedia
        Commons).
  \item Cap.~16: micrografía electrónica del panal de un convertidor catalítico, de
        Galadrid, CC BY-SA 4.0 (Wikimedia Commons).
  \item Cap.~17: rayos de sol y efecto Tyndall, de Kevin c higgins, CC BY-SA 4.0
        (Wikimedia Commons).
  \item Cap.~18: cristales de rubí y de esmeralda, de Robert M. Lavinsky, CC BY-SA 3.0
        (Wikimedia Commons).
  \item Cap.~20: cristales de ferroceno, de Andrea Sella, CC BY 4.0 (Wikimedia
        Commons).
  \item Cap.~22: un lingote de silicio monocristalino, de Sebastian Wallroth, dominio
        público (Wikimedia Commons).
  \item Cap.~23: una flor sobre aerogel de sílice, NASA/JPL, dominio público; la
        copa de Licurgo, de Marie-Lan Nguyen, CC BY 2.5 (Wikimedia Commons).
  \item Cap.~24: un cangrejo herradura del Atlántico, de Kaldari, CC0 (Wikimedia
        Commons).
  \item Cap.~27: flores de pelitre, de Soramimi, CC BY-SA 4.0 (Wikimedia
        Commons).
  \item Cap.~30: corteza de un tejo del Pacífico, de JOE BLOWE (theforestprimeval),
        CC BY-SA 2.0 (Wikimedia Commons).
  \item Cap.~32: una proliferación de algas vista por el satélite Landsat 8, USGS y NASA,
        dominio público.
  \item Cap.~33: un doble colector de vacío y de gas inerte, de Polimerek, CC BY-SA
        3.0 (Wikimedia Commons).
\end{itemize}
\clearpage
